\section{Calibration, rounding, and the posterior}
\label{sec:online}

We now use the sharp atom theorem to identify the minimum-cost scores
after each online history. This section supplies the complete argument;
no choice of a favorable minimizer is substituted for a posterior claim.

In this section $\mathcal C$ is the retained column-code space and
$\mathcal D=\mathcal D_{\mathcal A,\mathrm{odd}}$ is its retained
odd downward space from Section~\ref{sec:retained-stability}.
Let $P_{\mathcal C}F=b\,g(z)$, where $g(z)=\E[bF\mid z]$.
All projections below use the uniform physical input law.
A revealed logical fiber means that all its physical inputs have been
labeled. The target is positive relative to $b$ only at $z=0$.

Let $M=2^m$, $m\ge6$. Write
$w_s\ge w_0>0$ for code weights and
$W(z)=M^{-1}\sum_s w_s(-1)^{s\cdot z}$. Set $u=e_1+e_2$.
Take an affine coset $U=v_0+H$ with $K_U=|U|\le M/32$,
where $u\notin H$ and $v_0\notin H+\langle u\rangle$. Set
$H'=H+\langle u\rangle$, $U'=U+\langle u\rangle$, and
$E'=H'^\perp\setminus(H'+\langle v_0\rangle)^\perp$.
Assume $U'\cup\{u\}$ is independent in the Cayley graph whose
edges add one of the code columns $a_j$. Section~\ref{sec:path-code}
constructs such a coset containing exponentially many distinct costs.
Then $|E'|=M/(4K_U)\ge8$, $s\cdot u=0$ for $s\in E'$, and
\[
 \sum_{s\in E'}(-1)^{s\cdot z}
 =|E'|(\mathbf1_{H'}(z)-\mathbf1_{U'}(z)).
\]

For every nonzero coset $c$ of $\langle u\rangle$, put
\[
 k_c=\max\{0,W(z),W(z+u)\}\quad(z\in c),\qquad
 K_* =\sum_{c\ne0}k_c,\qquad w_*=W(u).
\]
The scalar bounds above imply, for all sufficiently large $m$,
\begin{equation}\label{eq:rect-dominance}
 0<w_*\le w_0/16,\qquad K_*\le w_*/4,\qquad M w_*\ge16w_0,
 \qquad W(v+u)>W(v)\ge0\quad(v\in U).
\end{equation}

\subsection{The common dual}

Put $L=M/2$ and define
\[
 h=\frac{(L-1)(w_*-4w_0/M)+K_*}{L-2},\qquad
 b_* =\frac{M(h+K_*)}{L-1},\qquad a_*=2w_0-b_*.
\]
These satisfy $3w_*/4\le h\le5w_*/4$, $b_*\le4w_*\le w_0/4$,
and $7w_0/4\le a_*\le2w_0$. Set
\[
 r_c=\frac{(L-1)k_c}{Lh+K_*},\qquad R_*=\sum_{c\ne0}r_c\le1/3.
\]
On $u^\perp$ define the strictly positive probability distribution
\[
 P(s)=\frac1L\left(1+\sum_{c\ne0}r_c(-1)^{s\cdot z_c}\right).
\]
Its mass at zero is
\[
 p=P(0)=\frac{h+K_*}{Lh+K_*},\qquad
 1-p=\frac{(L-1)h}{Lh+K_*}.
\]
Let $Q$ be $P$ conditioned on $s\ne0$, extended by zero outside
$u^\perp\setminus\{0\}$. Define a logical dual by
\begin{equation}\label{eq:rect-dual}
 \widehat\lambda(s)=w_s-2w_0\mathbf1_{s=0}
                    -MhQ(s)-\frac{a_*}{|E'|}\mathbf1_{s\in E'}.
\end{equation}
For every $s$, $|\widehat\lambda(s)|\le w_s$. Indeed its value at zero
is $-w_0$, and it equals $w_s$ if $s\cdot u=1$. Otherwise
$MhQ(s)\le Mh p/(1-p)=b_*$, so the total subtraction is at most
$b_*+a_*/|E'|\le2w_0\le2w_s$.

For $z\notin\{0,u\}$, Fourier orthogonality gives
\[
 \mathbb E_Q(-1)^{s\cdot z}
 =\frac{k_{[z]}}h-\frac{h+K_*}{(L-1)h}.
\]
Consequently
\begin{equation}\label{eq:rect-dual-values}
 \lambda(z)=M(W(z)-k_{[z]})-a_*
            -a_*\mathbf1_{H'}(z)+a_*\mathbf1_{U'}(z)
 \quad(z\ne0,u).
\end{equation}
The defining equation for $h$ also gives
\[
 \lambda(u)=0,\qquad \lambda(0)=M(W(0)-w_*)\ge a_*.
\]
In particular,
\[
 \lambda(v)=M\nu_v<0,\qquad \lambda(v+u)=0,
 \qquad \nu_v:=W(v)-W(v+u)\quad(v\in U).
\]
On the observed logical region
$\mathcal O=G\setminus(U'\cup\{u\})$, its sign is positive at zero
and negative elsewhere, with magnitude at least $a_*$.

There is also a useful strict coefficient slack. If $s\ne0$ and
$s\cdot u=0$, then
\begin{equation}\label{eq:rect-invariant-slack}
 \min\{w_s-\widehat\lambda(s),w_s+\widehat\lambda(s)\}\ge w_*.
\end{equation}
For the first term use
$MhQ(s)\ge2h(1-R_*)/(1-p)\ge w_*$.
For the second use
$2w_s-b_*-a_*/|E'|\ge a_*(1-1/|E'|)>w_*$.

\subsection{Observability outside a logical independent set}

\begin{lemma}\label{lem:independent-observability}
If $J\subseteq G$ is independent in the Cayley graph generated by the code columns $a_j$,
then every $A\in\mathcal D\cap\mathcal C^\perp$ satisfies
\[
 \|A\|_2^2\le2\mathbb E[\mathbf1_{z\notin J}A^2].
\]
\end{lemma}
\begin{proof}
Every Fourier support $T$ contributing to $A$ is an odd-degree proper
subset of a carrier-containing transversal. If it had a complete rail
at every pair, then either it would be a code support or, after losing
the carrier, would have even degree. Thus it has an incomplete rail
pair. Assign it to the first such pair $i$. Write $A=\sum_i A_i$ for this partition.

Conditioned on a logical fiber, two physical characters have nonzero
inner product only when their symmetric difference is a union of full
rail pairs with incidence vector $(a_j\cdot s)_j$ for some $s$. A symmetric difference containing the carrier has zero
conditional expectation as well, since the carrier remains an
independent uniform sign on each logical fiber. If one support is a proper subset of a rail at pair $i$,
its complement in that pair cannot be contained in a single rail.
Thus conditional covariance with another downward character requires
the same restriction at pair $i$. It follows that $A_i,A_j$ are
conditionally orthogonal whenever $i\ne j$, and that
$\ell_i(z):=\mathbb E[A_i^2\mid z]$ is invariant under $z\mapsto z+a_i$.
Indeed every surviving logical Fourier mode $s$ has $a_i\cdot s=0$.

An independent set contains at most one endpoint of each edge adding $a_i$.
Consequently
\[\mathbb E[\mathbf1_{z\in J}\ell_i(z)]\le\|A_i\|_2^2/2.\]
Sum over $i$, using conditional orthogonality, to obtain the assertion.
\end{proof}

By the assumed Cayley independent-set property,
Lemma~\ref{lem:independent-observability} applies to $\mathcal O$.

\subsection{Physical variance control}

Suppose a score $F$ and a nonnegative cost $B\le B_{\rm cap}$ satisfy
\[
 Q(F):=\sum_s w_s|\widehat F(A_s)|\le B,\qquad
 \|F\|_\infty\le H_{\rm cap},\qquad
 \|F-P_{\mathcal D}F\|_2^2\le S_{\rm at} B D,
\]
where $D=B-\langle b\lambda,F\rangle\ge0$ and $b$ is the physical
base character. These inequalities follow from Theorem~\ref{thm:atomic}. Its atomwise
conditions are
\[
 Q(\phi_a)\le\|a\|_2,\qquad
 \|\phi_a-P_{\mathcal D}\phi_a\|_2^2
 \le S_{\rm at}\|a\|_2(\|a\|_2-Q(\phi_a)),
 \quad \phi_a(x)=\clip(a\cdot\bar x),
\]
for a representation $F=\sum_jc_j\phi_{a_j}$ of cost
$B=\sum_j|c_j|\|a_j\|_2$. Indeed the dual inequality gives
$|\langle b\lambda,\phi_a\rangle|\le Q(\phi_a)$, and
Cauchy--Schwarz bounds the squared sum of residuals by $S_{\rm at}BD$.
The elementary bound $|\phi_a(x)|\le\sqrt{d+1}\|a\|_2$ also gives
$H_{\rm cap}=\sqrt{d+1}B_{\rm cap}$. For the retained column code, one may take
$S_{\rm at}=2^{20}\ell^2R^2$ by the subcode estimate in
Section~\ref{sec:retained-stability}.
Assume closed full-cube margins and correct labels outside $U$, with
logical target positive only at zero. Extra fully revealed fibers in
$U$ have the negative label. Write $P_{\mathcal C}F=bg$ and
\[
 S=\mathbb E[\mathbf1_{z\in\mathcal O}|\lambda(z)|
                         (fF-1/2)]\ge0.
\]
On each observed fiber, the variable $Y=fF$ is at least $1/2$.
Thus $\mathbb E[|Y-\mathbb EY|\mid z]\le
2(\mathbb E[Y\mid z]-1/2)$. Since $|\lambda|\ge a_*$ there,
this proves
$\mathbb E[\mathbf1_{\mathcal O}|F-P_{\mathcal C}F|]\le2S/a_*$.
Applying independent-set observability to
$P_{\mathcal D}F-P_{\mathcal C}F$ therefore yields
\begin{equation}\label{eq:rect-variance}
 L_F:=\|F-P_{\mathcal C}F\|_2^2
 \le C_{\rm phys}(D+S),\qquad
 C_{\rm phys}=\max\{16H_{\rm cap}/a_*,\ 5S_{\rm at}B_{\rm cap}\}.
\end{equation}
To see the constants, write $R=F-P_{\mathcal C}F=A+E$, where
$A=P_{\mathcal D}F-P_{\mathcal C}F$ and $E=F-P_{\mathcal D}F$.
The preceding conditional bound and $|R|\le2H_{\rm cap}$ give
$\mathbb E[\mathbf1_{\mathcal O}R^2]\le4H_{\rm cap}S/a_*$.
Observability and $(R-E)^2\le2R^2+2E^2$ imply
$\|A\|_2^2\le4\mathbb E[\mathbf1_{\mathcal O}R^2]+4\|E\|_2^2$.
Adding $\|E\|_2^2$ proves \eqref{eq:rect-variance}.

\subsection{Linear rounding to zero or a live rectangle}

Put $p_v=g(v)+1/2$ for $v\in U$ and
$\delta=D/w_*+S/a_*$. By \eqref{eq:rect-invariant-slack}, the nonconstant
$u$-invariant part $r$ of $g$ satisfies $\|r\|_\infty\le D/w_*$.
The region $\mathcal O_0=G\setminus(U'\cup\{0,u\})$ is invariant
under $u$, has density at least $1/2$, and consists of known negative
fibers. Averaging $g+1/2$ there eliminates the $u$-anti-invariant part.
Consequently
\[
 |\widehat g(0)+1/2|\le D/w_*+2S/a_*,\qquad
 \sup_z|g(z)+g(z+u)+1|\le4\delta.
\]
Each partner $v+u$ is a known negative fiber, so $p_v\ge-4\delta$.
In particular the total negative mass of $p$ is at most $M\delta$.

Compare $g$ with
\[
 g_{\rm ref}=-1/2+\mathbf1_0-\mathbf1_u
              +\sum_{v\in U}p_v(\mathbf1_v-\mathbf1_{v+u}).
\]
The logical $\ell^1$ norm of their difference is at most
\[
 2MS/a_*+4(K_U+1)\delta\le3M\delta.
\]
Here the observed fibers contribute at most $MS/a_*$; the fiber $u$
contributes at most the anchor excess plus $4\delta$; and the $K_U$
partner fibers contribute at most $4K_U\delta$.

There are characters $s$ with $s\cdot u=1$ and
$(-1)^{s\cdot v}=-1$ for every $v\in U$: prescribe the two independent linear conditions
$s\cdot u=1$ and $s\cdot v_0=1$ within $H^\perp$. For any such $s$, the dual is saturated on its
positive side, so $\widehat g(s)\ge-D/(2w_s)$. Since
\[
 \widehat g_{\rm ref}(s)=\frac2M\left(1-\sum_{v\in U}p_v\right),
\]
we obtain
\[
 \sum_{v\in U}p_v\le1+\frac{MD}{4w_0}+\frac32M\delta
 \le1+2M\delta.
\]
Thus the total positive mass is at most $1+3M\delta$.

On a fiber with $0<p_v<1$, the mean $g(v)=p_v-1/2$ lies
between the two allowed output intervals. A variable $Y=bF$ with
$|Y|\ge1/2$ and mean $g(v)$ has variance at least
$1/4-g(v)^2=p_v(1-p_v)$. Consequently
\[
 V_{\rm frac}:=\sum_{0<p_v<1}p_v(1-p_v)\le M L_F.
\]
Round each $p_v$ to its nearest $b_v\in\{0,1\}$. Negative coordinates
contribute at most $M\delta$ to the rounding distance. Coordinates above
one contribute at most $3M\delta$, because total positive mass is at most
$1+3M\delta$. The fractional coordinates contribute at most $2V_{\rm frac}$.
Hence $A:=\|p-b\|_1\le4M\delta+2V_{\rm frac}$.
If $b=0$, choose the zero vector. Otherwise choose one coordinate with
$b_v=1$; it is live because removed fibers have $p_v\le0$. The distance
from $b$ to this point mass is at most $3M\delta+A$. Therefore
\begin{equation}\label{eq:rect-rounding}
 \operatorname{dist}_1(p,\{0\}\cup\{e_v:v\text{ live}\})
 \le11M\delta+4V_{\rm frac}
 \le B_{\rm round}(D+S),
\end{equation}
where
\[
 B_{\rm round}=M(11/w_*+11/a_*+4C_{\rm phys}).
\]

\subsection{Absorption and classification of every minimum}

The baseline score $g_0=-1/2+\mathbf1_0-\mathbf1_u$ and rectangle scores
$g_v=g_0+\mathbf1_v-\mathbf1_{v+u}$ saturate the common dual. Their
weighted Fourier costs are
\[
 C_0=w_0/2+W(0)-W(u),\qquad C_v=C_0+\nu_v.
\]
The exact calibrated decomposition is
\[
 B=C_0+D+S+\sum_{v\in U}\nu_v p_v.
\]
No omitted-fiber term remains: $\lambda(u)=0$ and
$\lambda(v+u)=0$ for all $v\in U$.
Let $W_{\rm eff}=\max_{v\in U}|\nu_v|$. If
$W_{\rm eff}B_{\rm round}<1$, rounding gives
\[
 B\ge C_{\rm choice}+(1-W_{\rm eff}B_{\rm round})(D+S),
\]
where the choice is zero or a live point mass. The zero choice has cost
$C_0$, strictly greater than every live star because $\nu_v<0$.
Consequently any minimum-cost completion is a cheapest live rectangle.
Keep one live representative of each distinct effective cost, revealing
the remaining fibers during calibration. The live costs are pairwise
distinct, so the score is unique:
equality forces $D=S=0$, then $L_F=0$, exact anti-invariance around
$-1/2$, and the rounded point mass determine every logical and physical
value.

The path-code parameters in Section~\ref{sec:path-code} make every
constant other than $M=2^m$ in $B_{\rm round}$ polynomial in $m$.
That section constructs the affine set $U$ with
$W_{\rm eff}\le 2^{-3m}$ times any required fixed inverse polynomial,
by imposing a sufficiently small fixed fraction of path coordinates
with value one.
Hence $W_{\rm eff}B_{\rm round}<1$ for all sufficiently large $m$.
The strict pairwise cost ordering proved above identifies a unique
cheapest live rectangle after every history consisting of complete
revealed fibers.

\subsection{Every sufficiently large finite width}

\begin{lemma}[Integer allocation of a finite atom representation]
Let $F=\sum_{j=1}^J c_j\clip(a_j\cdot\bar x)$ be a finite
representation with nonzero $a_j,c_j$. Put
\[
 B=\sum_j|c_j|\|a_j\|_2,\qquad
 p_j=\frac{|c_j|}{\sqrt d\,\|a_j\|_2},\qquad p=\sum_jp_j.
\]
If $p<1$, then for every sufficiently large integer $N$ there is an
exact width-$N$ representation of $F$ with Gaussian energy tending to
$2\sqrt d\,B$ as $N\to\infty$.
\end{lemma}
\begin{proof}
Choose $r_j=\lceil Np_j\rceil$. For $N(1-p)\ge J$,
$\sum_jr_j\le N$. Use $r_j$ copies of the incoming vector $a_j$,
each with readout $c_j/r_j$, and pad with zero neurons. The output is
exactly $F$. Direct calculation gives
\[
 K_N-2\sqrt d\,B
 =\sum_j\frac{d\|a_j\|_2^2}{N}
                \frac{(r_j-Np_j)^2}{r_j}
 \le\frac dN\sum_j\|a_j\|_2^2.
\]
The right side tends to zero. The same statement applies after any
fixed small multiplicative increase of all readouts, provided the
resulting allocation sum remains less than one.
\end{proof}

The parity-orbit representation of $bg$ has cost $B_w(g)$ and
allocation sum
\[
 p=\frac1{\sqrt d}\sum_s
       \frac{|\widehat g(s)|}{\sqrt{n_s}\kappa_{n_s}}.
\]
Every baseline or rectangle has ordinary Fourier norm exactly $3/2$,
and $\sqrt n\kappa_n$ is bounded below by a positive absolute constant.
Thus $p=O(d^{-1/2})<1/2$ for large $m$. Integer allocation applies to
every witness and to a small strict-margin dilation of it.

Fix a history with at least one live fiber. Let $C_*$ be its unique
cheapest rectangle cost, and choose a polynomial cap $B_{\rm cap}$
strictly larger than all rectangle costs. The set of pairs $(B,F)$
satisfying the closed score inequalities above, $0\le B\le B_{\rm cap}$,
$\|F\|_\infty\le\sqrt{d+1}B_{\rm cap}$, the global closed margins,
and all observed labels is compact. The classification proof shows
that its unique minimum-cost score is the cheapest live rectangle.
At a physical input of that rectangle's exceptional fiber, intersect
this compact set with the condition of a correct prediction. The
resulting closed set has minimum cost $C_{\rm good}>C_*$.
If it is empty, define $C_{\rm good}=+\infty$.
Every representation of cost at most the cap satisfying the closed
full-cube and history margins satisfies these inequalities. Representations above the cap have larger energy by
\eqref{eq:variation-energy}. Therefore all feasible networks predicting correctly at the designated input have
energy at least
\[
 E_{\rm good}=2\sqrt d\,\min(C_{\rm good},B_{\rm cap})
              >2\sqrt d\,C_*.
\]
A small fixed dilation of the cheapest rectangle has strict margins
and variation cost below $E_{\rm good}/(2\sqrt d)$. Integer allocation
realizes it at every sufficiently large width with energy below
$E_{\rm good}$.

Fix such a width and let $P=N(d+2)$ be the parameter dimension.
A small open ball around this strict feasible wrong witness has energy
at most $a<E_{\rm good}$. Choose $\rho\in(0,1)$ with
$(1-\rho)E_{\rm good}>a$. The unnormalized Gaussian integral over
feasible networks predicting correctly at the designated input is at most
\[
 e^{-(1-\rho)E_{\rm good}/(2\sigma^2)}
 \int_{\R^P}e^{-\rho K_N/(2\sigma^2)}\,d\theta
 =C_{N,\rho}\sigma^P e^{-(1-\rho)E_{\rm good}/(2\sigma^2)}.
\]
The version-space integral is at least the ball volume times
$e^{-a/(2\sigma^2)}$. Their ratio tends to zero. This proves the
claimed actual posterior prediction at the designated input.

\subsection{The fixed nonrepeating sequence}

First reveal every physical input outside $U$, and every fiber in $U$
except one chosen representative of each distinct effective cost.
Sort the remaining live fibers by increasing rectangle cost. At each subsequent stage, take one fixed
physical input in the cheapest surviving fiber as the designated round,
then reveal the rest of that fiber before moving to the next one.
All labels are target labels, and no input is repeated. The sequence
depends only on the scalar costs, not on $N$ or $\sigma$.
There are at least $2^{(m-2)/3}=\exp(\Omega(m))$ designated rounds.
For each fixed construction there are finitely many stages; take the
maximum of their width thresholds. The posterior conclusion then
holds at every stage and every fixed width beyond that maximum.
Posterior-majority prediction is wrong once the correct-label
probability is below $1/2$, and posterior-sampling expected mistakes
converge to the designated round count by linearity of expectation.
